本节说明实验程序的实现思路。由于本实验的原理分析、攻击算法设计和运行结果都对应基础版程序，因此这里仅围绕 \verb|attack.py| 进行说明。

\subsection{程序结构}

\verb|attack.py| 的整体流程与第 3 节、第 5 节中的算法描述是一一对应的。程序先将目标密文拆成 $IV,C_1,C_2,C_3$ 四个分组，然后从最后一个目标分组开始，逐块恢复中间状态和明文。

代码中用到的几个辅助函数分别完成如下工作：
\begin{itemize}
    \item \verb|get_bytes|：从十六进制字符串中截取指定区间的字节；
    \item \verb|get_list|：将十六进制串按字节拆成列表；
    \item \verb|list_to_str|：将字节列表重新拼接成十六进制串。
\end{itemize}

这些辅助函数虽然比较简单，但使后面的分组截取、字节替换和伪造密文拼接都更加直观，也更接近实验推导时的写法。

\subsection{主要实现流程}

基础版程序的外层循环按照
\[
    C_3 \rightarrow C_2 \rightarrow C_1
\]
的顺序逐块恢复明文。对每个目标分组，程序的主要步骤如下：

\begin{enumerate}
    \item 取当前目标密文分组 $y$ 及其真实前一密文分组 $r$；
    \item 先通过修改 $r$ 的不同位置，判断原始填充长度；
    \item 根据原始填充长度，先恢复中间状态末尾已经确定的若干字节；
    \item 再不断调整伪造分组后缀，并枚举新的字节，逐步向前恢复整个中间状态；
    \item 最后将恢复出的中间状态与真实前一密文分组异或，得到当前明文分组。
\end{enumerate}

从实现方式上看，这个程序并没有把所有步骤压成过于抽象的函数，而是比较直接地按照攻击过程写出来，因此和实验指导书中的推导过程对应得比较紧，也更能体现“先判断填充长度，再继续向前恢复”的思路。

\subsection{实现体会}

从实际编写过程来看，Padding Oracle 攻击真正麻烦的地方不在公式，而在于如何把“逐字节构造、逐次提交、根据反馈修正”这些步骤准确地落到代码里。尤其是在基础版实现中，字节位置、填充值和中间状态之间的关系必须完全对应，否则后面恢复出来的结果就会整体出错。

虽然基础版程序运行速度较慢，但它把实验指导书中的思路保留得最完整，也最适合作为本实验报告正文的实现说明。
